\section{Scope and further questions}\label{sec:discussion}

The continuous minimax formulas, finite-grid characterizations, and instance
algorithms address different decisions. A minimax formula specifies the worst
rounding error before the relaxed input is known. An exact instance algorithm
finds the best schedule after the input is supplied. A coarsening certificate
bounds the loss caused by reducing the allowed switching grid. Their
combination separates switch-budget error, grid restriction, and input
representation error without imposing smoothness on the original measurable
control in the minimax analysis.

Three mathematical extensions remain substantial. First, the exact
five-block one-sided reach bound for general mode count is unresolved. The
exclusion identity in Appendix~\ref{sec:higher-reach} identifies the missing
weighted three-block inequality. The negative event assignment is not a
control trajectory; chronological monotonicity removes that particular
obstruction. A certificate for one chamber does not cover the remaining
orderings or prove the full reach bound. The general and seeded results in
Section~\ref{sec:general-budgets} remain valid independently of this question.
A successful extension would improve higher-budget minimax formulas and the
second-order mode-count gap.

Second, the regime with as many or more activation blocks than modes has a
different structure. The three-mode one-sided obstruction rules out a direct
extension of the distinct-mode method. The exact full value $F_{3,2}(T)=T/5$
and the range $T/7\le F_{3,3}(T)\le T/6$ show that finite-grid chamber methods
can nevertheless improve the boundary theory. The floor-history construction
is valid for every grid length, but the enumerations here stop at seven cells.
The six exceptional seven-cell chambers have explicit repairs; their existence
also disproves the tempting unit-error extension on $2s+1$ cells. Determining
$F_{3,3}$ and a general three-mode floor-history switch law requires further
arguments. For other small-mode regimes, the classical bound
\eqref{eq:classical-small-mode} remains a useful comparison.

Third, constraints beyond a switch budget change the approximation question.
The exact subset algorithm handles its stated per-mode minimum dwell rules,
but a transition graph couples successive labels and is outside that unary
assignment reduction. The odd-grid dwell example shows that unrestricted
coarsening certificates can fail even under arbitrarily fine refinement.
Useful constrained certificates would need assumptions ensuring feasible
nearby switching times, or a constraint-aware refinement construction.

The algorithmic results already support exact small-budget rounding and
certified additive approximation with rational data. A broader computational
study could compare matched budgets and objectives against established graph
and mixed integer methods, and then evaluate resulting trajectories under
specified dynamics. Such a study would address application performance;
it would not change the scope of the discrepancy theorems.
